% pretextual.tex — elementos pré-textuais (capa, folha de rosto, resumos, listas, sumário)
% Dados fornecidos pelo autor em 2026-09-10: campus, cidade e orientação.
% [VALIDAR COM O AUTOR]: curso exato, banca e demais elementos institucionais.

\frontmatter

% ---------------------------------------------------------------------------
% Capa
% ---------------------------------------------------------------------------
\begin{center}
  \textbf{INSTITUTO FEDERAL FLUMINENSE}\\[2pt]
  Campus: Meu laboratório pessoal\\
  Curso de Engenharia de Controle e Automação [VALIDAR COM O AUTOR]
\end{center}

\vspace{2.5cm}

\begin{center}
  \textbf{WILLIAM DA SILVA VIANNA}
\end{center}

\vspace{3.0cm}

\begin{center}
  \textbf{LABORATÓRIO 404: DESENVOLVIMENTO DE UM JOGO SÉRIO 3D COM\\
  INTELIGÊNCIA ARTIFICIAL LOCAL PARA APOIO AO\\ ENSINO DE AUTOMAÇÃO INDUSTRIAL}
\end{center}

\vfill

\begin{center}
  Campos dos Goytacazes\\
  2026
\end{center}

\newpage

% ---------------------------------------------------------------------------
% Folha de rosto
% ---------------------------------------------------------------------------
\begin{center}
  \textbf{WILLIAM DA SILVA VIANNA}
\end{center}

\vspace{2.5cm}

\begin{center}
  \textbf{LABORATÓRIO 404: DESENVOLVIMENTO DE UM JOGO SÉRIO 3D COM\\
  INTELIGÊNCIA ARTIFICIAL LOCAL PARA APOIO AO\\ ENSINO DE AUTOMAÇÃO INDUSTRIAL}
\end{center}

\vspace{2.5cm}

\hfill
\begin{minipage}{8cm}
  \SingleSpacing
  \footnotesize
  Monografia apresentada ao curso de Engenharia de Controle e Automação
  [VALIDAR COM O AUTOR] do Instituto Federal Fluminense, como requisito
  parcial para obtenção do título de bacharel.\\
  Orientador: o próprio autor.\\
  Área de concentração: automação industrial e tecnologia educacional.
\end{minipage}

\vfill

\begin{center}
  Campos dos Goytacazes\\
  2026
\end{center}

\newpage

% ---------------------------------------------------------------------------
% Resumo
% ---------------------------------------------------------------------------
\chapter*{Resumo}
\addcontentsline{toc}{chapter}{Resumo}

\SingleSpacing

Este trabalho descreve o projeto, o desenvolvimento e a validação técnica do
Laboratório 404, um jogo sério 3D educacional/ficcional no qual uma
inteligência artificial local --- a personagem \aria{} --- participa do mundo
por meio de diálogo, sem nunca controlar o jogo. O problema abordado é
triplo: integrar um grande modelo de linguagem (LLM) a um jogo 3D mantendo
segurança e previsibilidade; executar essa integração em um computador pessoal
sem GPU dedicada; e conduzir o desenvolvimento de forma rastreável, com
requisitos verificáveis e evidências reproduzíveis. A solução foi construída em
seis provas de conceito incrementais (\poc{1} a \poc{6}) sobre o motor Godot~4.5,
com assets 3D modelados no Blender~5.2.1 e um adaptador local em Python/FastAPI
que media a conversa com o Ollama, restringindo a saída do modelo a cinco
intenções validadas em duas camadas independentes, com resposta de
\emph{fallback} obrigatória. A validação combinou uma especificação com
requisitos funcionais e não funcionais, critérios de aceitação no formato
DADO/QUANDO/ENTÃO, uma suíte automatizada \emph{headless} e evidências visuais
geradas por um arnês de captura. Os resultados obtidos na máquina de referência
(AMD Ryzen~7 5700U, GPU integrada) foram: 211 verificações do motor e 16 testes
do adaptador aprovados; cena 3D com 342 nós, 37 materiais e 1.018,7~KB em GLB,
executando a 60~FPS limitados por sincronismo vertical em 1280$\times$720; e
latências da conversa com o modelo local \texttt{llama3.1:8b} entre 12,9~s e
22,0~s em três contextos medidos. Conclui-se que a integração proposta é viável
no nível de protótipo, com limites declarados: ausência de estudo com usuários,
validação em hardware real pendente, modelo-alvo \texttt{llama3.2:3b} não
executado e empacotamento independente ainda não gerado.

\vspace{1em}
\noindent\textbf{Palavras-chave:} jogos sérios; automação industrial; grandes
modelos de linguagem; inferência local; Godot; desenvolvimento orientado a
especificação.

\newpage

% ---------------------------------------------------------------------------
% Abstract
% ---------------------------------------------------------------------------
\chapter*{Abstract}
\addcontentsline{toc}{chapter}{Abstract}

\SingleSpacing

This work describes the design, development and technical validation of
Laboratório 404, a serious educational/fictional 3D game in which a local
artificial intelligence --- the character \aria{} --- participates in the world
through dialogue, without ever controlling the game. The problem addressed is
threefold: integrating a large language model (LLM) into a 3D game while
preserving safety and predictability; running this integration on a personal
computer without a dedicated GPU; and conducting the development in a traceable
way, with verifiable requirements and reproducible evidence. The solution was
built in six incremental proofs of concept (\poc{1} to \poc{6}) on the Godot~4.5
engine, with 3D assets modeled in Blender~5.2.1 and a local Python/FastAPI
adapter that mediates the conversation with Ollama, restricting the model output
to five intents validated in two independent layers, with a mandatory fallback
response. Validation combined a specification with functional and non-functional
requirements, acceptance criteria in GIVEN/WHEN/THEN format, an automated
headless test suite and visual evidence produced by a screenshot harness. Results
on the reference machine (AMD Ryzen~7 5700U, integrated GPU) were: 211 engine
checks and 16 adapter tests passing; a 3D scene with 342 nodes, 37 materials and
1018.7~KB in GLB, running at 60~FPS vsync-limited at 1280$\times$720; and
conversation latencies with the local \texttt{llama3.1:8b} model between 12.9~s
and 22.0~s across three measured contexts. It is concluded that the proposed
integration is feasible at prototype level, with declared limitations: no user
study, pending validation on real hardware, target model \texttt{llama3.2:3b}
not executed, and standalone packaging not yet produced.

\vspace{1em}
\noindent\textbf{Keywords:} serious games; industrial automation; large language
models; local inference; Godot; specification-driven development.

\newpage

% ---------------------------------------------------------------------------
% Listas
% ---------------------------------------------------------------------------
\SingleSpacing

\listoffigures
\newpage
\listoftables
\newpage
\listof{quadro}{Lista de Quadros}
\newpage

\chapter*{Lista de Abreviaturas e Siglas}
\addcontentsline{toc}{chapter}{Lista de Abreviaturas e Siglas}
\noindent
\begin{tabular}{@{}>{\bfseries}p{2.2cm} p{12.5cm}@{}}
ABNT & Associação Brasileira de Normas Técnicas \\
API & \emph{Application Programming Interface} (interface de programação de aplicações) \\
CLP & Controlador lógico programável \\
CPU & \emph{Central Processing Unit} (unidade central de processamento) \\
FPS & \emph{Frames per second} (quadros por segundo) \\
GLB & \emph{glTF Binary} (formato binário de glTF) \\
glTF & \emph{Graphics Language Transmission Format} \\
GPU & \emph{Graphics Processing Unit} (unidade de processamento gráfico) \\
HUD & \emph{Head-Up Display} (interface de informações em tela) \\
IA & Inteligência artificial \\
IFF & Instituto Federal Fluminense \\
JSON & \emph{JavaScript Object Notation} \\
LLM & \emph{Large Language Model} (grande modelo de linguagem) \\
MCP & \emph{Model Context Protocol} \\
NPC & \emph{Non-Player Character} (personagem não jogável) \\
POC & \emph{Proof of Concept} (prova de conceito) \\
REST & \emph{Representational State Transfer} \\
SDD & \emph{Specification-Driven Development} (desenvolvimento orientado a especificação) \\
SO & Sistema operacional \\
WAV & \emph{Waveform Audio File Format} \\
\end{tabular}

\newpage

% ---------------------------------------------------------------------------
% Sumário
% ---------------------------------------------------------------------------
\SingleSpacing
\tableofcontents

\cleardoublepage
